% Figure 2.x (fig:bg:solver-path): the computational path of a classical PDE solver.
% Drawn in TikZ (loaded by jkureport.sty); colours from style/thesis-colors.sty.
% Message: a stencil turns the grid values into a rate f_i, a time step turns the
% rate into the next grid, and repeating the step is the rollout. One-dimensional
% schematic; the cell values are illustrative, not data. Explanations are in the
% caption (figure rule, NOTES 2026-10-18).
\begin{tikzpicture}[
    font=\small,
    >=latex,
    fixed/.style={draw=gaPrimary, fill=gaPrimaryTint, rounded corners=2pt,
                  minimum height=11mm, inner sep=4pt, align=center},
    cell/.style={draw=gaLine, minimum size=3.4mm, inner sep=0pt, line width=0.4pt},
    flow/.style={->, draw=gaInk, line width=0.6pt},
    loop/.style={->, draw=gaInkMuted, line width=0.6pt, dashed},
  ]
  \def\cw{0.34}
  % --- continuous field u(t_k, x) --------------------------------------------
  \draw[gaInkMuted, line width=0.4pt, ->] (0, 0) -- (2.0, 0) node[right, inner sep=1pt, text=gaInkMuted] {$x$};
  \draw[gaInkMuted, line width=0.4pt, ->] (0, 0) -- (0, 1.25);
  \draw[gaPrimary, line width=1.0pt, domain=0.05:1.9, samples=60, smooth]
    plot (\x, {0.12 + 0.85*exp(-((\x - 1.0)^2)/0.18)});
  % the seven cell values of u^k sample the field at the cell centres
  \foreach \j in {0, ..., 6} {
    \pgfmathsetmacro{\xs}{0.25 + 0.25*\j}
    \fill[gaPrimary] (\xs, {0.12 + 0.85*exp(-((\xs - 1.0)^2)/0.18)}) circle (1.1pt);
  }
  \node[text=gaInk] at (1.0, -0.42) {$u(t_k, x)$};

  % --- grid at step k: cell i highlighted together with its neighbours ---------
  \begin{scope}[shift={(2.75, 0.45)}]
    \foreach \v [count=\j from 0] in {11, 28, 65, 88, 65, 28, 11} {
      \node[cell, fill=gaPrimary!\v] (a\j) at (\j*\cw, 0) {};
    }
    \draw[gaPrimary, line width=1.0pt] ([xshift=-0.2pt]a2.north west) rectangle ([xshift=0.2pt]a4.south east);
    \draw[gaAccent, line width=1.2pt] (a3.north west) rectangle (a3.south east);
    \node[text=gaInk] at (3*\cw, -0.85) {$u^{k}$};
    \draw[gaInkMuted, line width=0.5pt]
      ([yshift=-1.2mm]a2.south west) -- ++(0, -0.6mm) -| ([yshift=-1.2mm]a4.south east);
    \node[text=gaInkMuted, font=\footnotesize] at (3*\cw, -0.47) {stencil};
  \end{scope}

  % --- time step -----------------------------------------------------------------
  \node[fixed] (step) at (7.55, 0.45) {time step\\$u_i^{k+1} = u_i^{k} + \Delta t\, f_i$};

  % --- grid at step k + 1 ----------------------------------------------------------
  \begin{scope}[shift={(10.05, 0.45)}]
    \foreach \v [count=\j from 0] in {12, 34, 70, 84, 70, 34, 12} {
      \node[cell, fill=gaPrimary!\v] (b\j) at (\j*\cw, 0) {};
    }
    \draw[gaAccent, line width=1.2pt] (b3.north west) rectangle (b3.south east);
    \node[text=gaInk] at (3*\cw, -0.87) {$u^{k+1}$};
  \end{scope}

  % --- arrows ------------------------------------------------------------------------
  \draw[flow] (2.15, 0.45) -- ([xshift=-1.5mm]a0.west);
  \draw[flow] ([xshift=1.5mm]a6.east) -- node[above, text=gaInk] {$f_i$} (step.west);
  \draw[flow] (step.east) -- ([xshift=-1.5mm]b0.west);

  % --- rollout: the new grid is the input of the next step --------------------------
  \draw[loop] ([yshift=1mm]b3.north) -- ++(0, 0.6)
    -| node[pos=0.25, above, text=gaInkMuted] {rollout} ([yshift=1mm]a3.north);
\end{tikzpicture}
